\cleardoublepage
\chapter*{Conclusions}\label{ch:conclusions}
\addcontentsline{toc}{chapter}{Conclusions}
The first part of this thesis explored the interaction mechanisms of charged particles with matter focusing on their implications in the fields of Particle Therapy (PT) and Space Radiation Protection (SRP), laying the groundwork for the study presented in this thesis. In particular, it was shown that nuclear interactions can significantly affect both of these applications, although they occur with a lower probability when compared to electromagnetic interactions. Indeed, nuclear processes mainly lead to the fragmentation of both the target and the projectile nuclei, producing secondary particles that may have undesirable effects not only in PT, impacting the dose distribution and the effectiveness of the treatment, but also in the context of SRP, causing radiation-induced damages to both astronauts and electronic systems employed in spacecraft for long-term space missions.

In order to investigate the nuclear fragmentation mechanisms more thoroughly, the FOOT (FragmentatiOn of Target) experiment was devised to measure the nuclear fragmentation double-differential cross sections (in the fragments' kinetic energy and emission angle) of ions with charge $Z\leq8$ and kinetic energies $100$--$800\,\mathrm{MeV/u}$, filling the gap of lacking experimental measurements relevant for both PT and SRP. The FOOT experiment will provide new measurements of both target and projectile fragmentation cross sections with a $5$--$10\%$ precision, establishing more stringent experimental constraints on theoretical models, deterministic codes and Monte Carlo simulations employed in PT and SRP to describe nuclear interactions.

To reach the requested FOOT accuracy on cross section measurements, a high acquisition rate and a precise reconstruction of particles' trajectories and momenta are required. This thesis work contributed to these aspects by implementing on the hardware side the trigger firmware for the new Vertex Tracker (VT), whose operation is foreseen as of November 2026, and by optimizing on the software side the global tracking system, enhancing the reconstruction of the current VT and the Microstrip Silicon Detector (MSD).

As to the hardware work, the upcoming data taking phase of the FOOT experiment demands an increase in data acquisition rate capabilities, necessary to raise the statistics for more accurate cross section measurements. A primary bottleneck in the original FOOT electronic setup was the readout time of the MIMOSA-28 pixel sensors in the current VT, which operated on a rolling-shutter mechanism requiring $185.6\,\mu\mathrm{s}$ per frame and restraining the DAQ rate to approximately $1\,\mathrm{kHz}$. The planned upgrade of the MIMOSIS-2.1 VT, with a global-shutter readout of $5\,\mu\mathrm{s}$, will remove this bottleneck. To support the new TDAQ architecture for the MIMOSIS-2.1 VT, this thesis work optimized the \texttt{VHDL} firmware of the VT trigger board (the TrgBox), which consists of a Terasic DE10-Nano board housing an Intel/Altera System-on-Chip FPGA (Cyclone V) and a dual-core ARM Cortex-A9 CPU. Specifically, the enhancement of the FPGA system introduced the configurable show-ahead (FWFT) FIFO in the event readout, the improvement of the post-trigger busy activation and the completion of the event-building logic synchronization. In addition, the \texttt{C++} software running on the HPS to configure the registers, monitor the trigger counters and send the events to the central DAQ was developed. The TrgBox operation was validated with an FPGA-based detector emulator and at the BTF campaign in Frascati laboratories, showing stable clock distribution, no event loss and reliable operation up to a trigger rate of $20\,\mathrm{kHz}$, proving ready for the CNAO data taking scheduled for November 2026.

In parallel with the hardware development, the offline tracking algorithm of the GENFIT-SHOE framework was optimized for both the current MIMOSA-28 VT and the MSD, using experimental data collected during the CNAO 2025 campaign, which studies nuclear fragmentation irradiating $200\,\mathrm{MeV/u}$ carbon ions on a $5\,\mathrm{mm}$ carbon target.

Regarding the VT updates, the pull distributions, defined as the ratio between the position residuals and their corresponding uncertainties, were investigated for the MIMOSA-28 sensors. They revealed non-zero means ($\mu_\text{pull}$) in almost all planes, and pull widths ($\sigma_\text{pull}$) up to about $5$ for the lightest fragments. First, a $\chi^2$ minimization algorithm was developed to estimate the longitudinal displacement $\Delta z$ of the VT planes by analyzing track slopes and residuals. This method pointed to a displacement of VT1, later confirmed by direct measurements of sensor distances. The corrected geometry reduced the estimated $\Delta z$ from $323\,\mu\mathrm{m}$ to $75.6\,\mu\mathrm{m}$ (extrapolated from the VT1 helium fit), validating the misalignment suggested by the non-zero pull means. This correction mostly improved the $\sigma_\text{pull}$ of light, most deflected fragments (for helium in VT0 and VT1 from $3.6$--$5.3$ to $\simeq2$). After the application of the software-based alignment procedure, the $\sigma_\text{pull}$ of carbons are found at around $1$, the reference value, whereas the lower-charge fragments they slightly remain above unity ($1.4$--$2.2$ for helium), indicating that the VT uncertainty model still underestimates the errors of narrower clusters. The inclusion of the covariance term in the cluster error and the study of residuals as a function of the cluster position (suggesting rotations of the planes of about $\pm1\,\mathrm{mrad}$ around the beam axis, not yet handled by the software alignment) indicate the direction for further refinements. As to the pull means, they were brought towards the reference value $0$ for all fragments only after applying the software-based alignment. In particular, the range of the carbon $\mu_{\text{pull}}$ was remarkably reduced from $[-0.6,1]$ to $[-0.1,0.1]$.

Regarding the MSD optimization, the original center-of-gravity clustering algorithm of the MSD overestimates the position uncertainties, with pull widths of about $0.3$--$0.6$ for all fragments after the VT1 geometry correction and the software-based alignment. To unravel this issue, a non-linear position finding algorithm based on the $f(\eta)$ integral method was implemented in SHOE, exploiting the charge sharing between adjacent readout strips. The $f(\eta)$ function was parameterized by a linear term plus four error functions calibrated on carbon ion data from all $6$ MSD planes, and successively produced a uniform spatial position mapping across the $150\,\mu\text{m}$ strip pitch. Employing the novel $f(\eta)$ uncertainty estimator brought the MSD $\sigma_{\text{pull}}$ to about $0.7$--$0.8$ for carbon and closer to $1$ for lighter fragments, improving considerably the overall MSD error estimation. Furthermore, a detailed analysis of the cluster signal showed that the few large position uncertainties assigned by $f(\eta)$ stem from low-signal noisy clusters, and, most importantly, that the MSD signal can better identify out-of-target fragmentation occurring in the air before the ToF Wall detector.

The last part of this thesis evaluated the reconstruction performances studying the VT and MSD improvements on the GENFIT-based global tracking algorithm and, successively, resolution on particle momentum. For carbon tracks, the $f(\eta)$ algorithm strongly suppressed the tail at higher $\chi^2$ values, shifting the mean of the reduced $\chi^2$ distribution from $1.54$ to $1.40$ (closer to the expected $1$) and reducing its standard deviation from $1.69$ to $1.01$. This variation suggests that the position and uncertainty model based on $f(\eta)$ partly discards badly-reconstructed tracks. In addition, the $f(\eta)$ significantly reduced the rise at large $p$-values caused by the MSD overestimated errors, rendering the $p$-value distribution closer to a uniform one. On the other hand, the momentum resolution was determined on experimental data by comparing the GENFIT fitted momentum with a quasi-independent momentum estimate from the time of flight and subtracting the latter contribution in quadrature. The final resolution resulted in about $3.34\%$ for \ce{^{12}C} primaries, in agreement with the $\simeq 3\%$ of Monte Carlo simulations, and about $2.96\%$ for \ce{^{4}He} fragments. Both values are compatible with the $3$--$5\%$ precision in momentum reconstruction required by the FOOT electronic setup.

In conclusion, the implementation of the TrgBox and the optimization of the tracking system will strengthen the reconstruction capabilities of the FOOT experiment, providing a solid foundation for its future experimental program and forthcoming campaigns at CNAO and GSI facilities.